% Generated draft: requires mathematical and source review.
\documentclass[11pt]{article}
\usepackage[margin=25mm]{geometry}
\usepackage{fontspec}
\setmainfont{Latin Modern Roman}
\usepackage{amsmath,amssymb,mathtools,unicode-math}
\setmathfont{Latin Modern Math}
\usepackage{enumitem,needspace,xurl,hyperref}
\setlength{\parindent}{0pt}
\setlength{\parskip}{5pt}
\setlength{\emergencystretch}{3em}
\newcommand{\N}{\mathbb{N}}
\newcommand{\Z}{\mathbb{Z}}
\newcommand{\Q}{\mathbb{Q}}
\newcommand{\R}{\mathbb{R}}
\newcommand{\C}{\mathbb{C}}
\newcommand{\F}{\mathbb{F}}
\newcommand{\A}{\mathbb{A}}
\newcommand{\PP}{\mathbb{P}}
\newcommand{\E}{\mathbb{E}}
\newcommand{\eps}{\varepsilon}
\newcommand{\dd}{\,\mathrm{d}}
\begin{document}
\section*{UnsolvedMath ID: 8800037}
\textbf{Source problem code:} catalogue entry
\subsection{Definitions and mathematical statement}
For an elliptic curve \(E/\F_q\), let \(N_q(E)=\#E(\F_q)\), including the identity point. It lies in the Hasse interval
\[
 q+1-2\sqrt q\leq N_q(E)\leq q+1+2\sqrt q.
\]
Its arithmetic structure concerns the integer factorization
\[
 N_q(E)=\prod_{\ell\text{ prime}}\ell^{v_\ell(N_q(E))},
\]
for example its prime factors, valuations, smoothness, or primality, rather than merely its magnitude or the abstract group structure of \(E(\F_q)\).

Shparlinski asks to characterize the \emph{typical} arithmetic structure of these orders as elliptic curves over finite fields vary. A precise distribution statement must specify how curves and fields are sampled and which arithmetic statistics and limiting regime are intended. Uniform Weierstrass equations, uniform isomorphism classes, and classes weighted by reciprocal automorphism-group size are different sampling rules. The source also does not say whether \(q\) runs over primes, all prime powers, or powers of a fixed characteristic.

The supplied summary discusses order distribution in the Hasse interval and partial factorization results. Its informal near-uniform distribution claim does not supply a precise measure, estimate, or factorization law. The present task is not reduced to the Hasse size bound, and no numerical threshold for smoothness or a particular limiting distribution is invented.
\subsection{Short English statement}
Describe how the prime factorization of an elliptic curve's finite-field order typically behaves. The source leaves the sampling of curves and fields, arithmetic statistics, and asymptotic regime unspecified.
\subsection{Sources}
\begin{itemize}
\item Local UnsolvedMath registry, permanent ID 8800037, source code catalogue entry; attributed to Igor Shparlinski; no proposal year supplied; update 2026-08-06.
\item IPAM, \emph{Workshop on Number Theory and Cryptography --- Open Problems} (2006), attachment 6224, PDF page 3: \url{https://www.ipam.ucla.edu/programs/workshops/workshop-i-number-theory-and-cryptography-open-problems/?tab=schedule}. Registry access date 2026-07-29; PDF and legacy-PPT visual audit. Provenance reports 19 surviving attachments, 17 PDFs and two legacy PPT decks, with repetitions and solved definitions excluded. Present status and redistribution terms require review.
\item The summary names Lenstra (1987), Deuring (1941), and Luca--Shparlinski for related distribution and arithmetic results without full bibliographic entries or exact hypotheses. Its reference to ``0041'' for smooth orders is a supplied cross-reference, not a separate verified theorem.
\end{itemize}
\end{document}
